% Need to use lualatex not xelatex
\documentclass[a4paper,10pt]{extarticle}
\usepackage[margin=0.1in,foot=0in,landscape]{geometry}
\input{LaTeX-ToolKit/template.tex}
\usepackage{multicol}
\lstset{
basicstyle=\ttfamily\fontsize{4pt}{4pt}\selectfont,
keywordstyle=\color{blue},
stringstyle=\color{red},
commentstyle=\color{green!50!black},
numbers=left,
stepnumber=0,
backgroundcolor=\color{white},
showspaces=false,
showstringspaces=false,
showtabs=false,
frame=none,
rulecolor=\color{black},
tabsize=4,
captionpos=b,
breaklines=true,
breakatwhitespace=false
}
\title{NTUEE Data Structure 115 Midterm Cheat Sheet}
\author{沈威宇}
\date{\temtoday}
\pagestyle{empty}
\begin{document}
\temdoc
\fontsize{5pt}{6pt}\selectfont
\begin{multicols}{7}

\tb{Definition}

For real-valued nets $f$ and $g$,

$\exists M>0,T\in\bbR\colon x\ge T\implies\abs{f(x)}\le M\abs{g(x)}$ iff $f(x)\in O(g(x))$ and $\abs{g(x)}$ is an asymptotic upper bound of $\abs{f(x)}$.

$\exists M>0,T\in\bbR\colon x\ge T\implies\abs{f(x)}\ge M\abs{g(x)}$ iff $f(x)\in \Omega(g(x))$ and $\abs{g(x)}$ is an asymptotic lower bound of $\abs{f(x)}$.

$\exists M_1,M_2>0,T\in\bbR\colon x\ge T\implies M_1\abs{g(x)}\le\abs{f(x)}\le M_2\abs{g(x)}$ iff $f(x)\in \Theta(g(x))$.

$\forall M>0\colon\exists T\in\bbR\colon x\ge T\implies\abs{f(x)}<M\abs{g(x)}$ iff $f(x)\in o(g(x))$.

$\forall M>0\colon\exists T\in\bbR\colon x\ge T\implies\abs{f(x)}>M\abs{g(x)}$ iff $f(x)\in \omega(g(x))$.

$\in$ is also denoted as $=$ for them.

\tb{Limit}

If $\exists k\in\bbR\colon x\ge k\implies g(x)\ne0$,

$\limsup\abs{\frac{f(x)}{g(x)}}<\infty$ iff $f(x)\in O(g(x))$.

Proof:
\(\abs{f(x)}\le M\abs{g(x)}\qc x\ge T\); 
\(\abs{\frac{f(x)}{g(x)}}\le M\qc x\ge T\); 
\(\sup\abs{\frac{f(x)}{g(x)}}\le M\qc x\ge T\); 
\(\limsup\abs{\frac{f(x)}{g(x)}}\le M<\infty\); 


$\liminf\abs{\frac{f(x)}{g(x)}}>0$ iff $f(x)\in \Omega(g(x))$.

Proof:
\(\abs{f(x)}\ge M\abs{g(x)}\qc x\ge T\); 
\(\abs{\frac{f(x)}{g(x)}}\ge M\qc x\ge T\); 
\(\inf\abs{\frac{f(x)}{g(x)}}\ge M\qc x\ge T\); 
\(\liminf\abs{\frac{f(x)}{g(x)}}\ge M>0\); 


$\infty>\limsup\abs{\frac{f(x)}{g(x)}}\ge\liminf\abs{\frac{f(x)}{g(x)}}>0$ iff $f(x)\in \Theta(g(x))$.

Proof:
\(\abs{f(x)}\ge M_1\abs{g(x)}\qc x\ge T\); 
\(\abs{\frac{f(x)}{g(x)}}\ge M_1\qc x\ge T\); 
\(\inf\abs{\frac{f(x)}{g(x)}}\ge M_1\qc x\ge T\); 
\(\liminf\abs{\frac{f(x)}{g(x)}}\ge M_1>0\); 
\(\abs{f(x)}\le M_2\abs{g(x)}\qc x\ge T\); 
\(\abs{\frac{f(x)}{g(x)}}\le M_2\qc x\ge T\); 
\(\sup\abs{\frac{f(x)}{g(x)}}\le M_2\qc x\ge T\); 
\(\limsup\abs{\frac{f(x)}{g(x)}}\le M_2<\infty\); 


$\lim_{x\to\infty}\abs{\frac{f(x)}{g(x)}}=0$ iff $f(x)\in o(g(x))$.
\(\abs{f(x)}<M\abs{g(x)}\); 
\(\abs{\frac{f(x)}{g(x)}}<M\); 
\(\lim_{x\to\infty}\abs{\frac{f(x)}{g(x)}}=0\); 

$\lim_{x\to\infty}\abs{\frac{f(x)}{g(x)}}=\infty$ iff $f(x)\in \omega(g(x))$.
\(\abs{f(x)}>M\abs{g(x)}\); 
\(\abs{\frac{f(x)}{g(x)}}>M\); 
\(\lim_{x\to\infty}\abs{\frac{f(x)}{g(x)}}=\infty\); 

\tb{Relation}

For any fixed $g(x)$, $o(g(x))$, $\Theta(g(x))$, and $\omega(g(x))$ are mutually exclusive.

Proof:
Prove by contradiction. If $f\in o(g(x))\cap\Theta(g(x))$,
\(\forall M>0\colon\exists T\in\bbR\colon x\ge T\implies\abs{f(x)}<M\abs{g(x)}\); 
\(\exists M_1,M_2>0,U\in\bbR\colon x\ge U\implies M_1\abs{g(x)}\le\abs{f(x)}\le M_2\abs{g(x)}\); 
then set $M=M_1$, for $x\ge\max\{T,U\}$
\(M_1\abs{g(x)}\le\abs{f(x)}<M_1\abs{g(x)},\); 
which can't happen. If $f\in \omega(g(x))\cap\Theta(g(x))$,
\(\forall M>0\colon\exists T\in\bbR\colon x\ge T\implies\abs{f(x)}>M\abs{g(x)}\); 
\(\exists M_1,M_2>0,U\in\bbR\colon x\ge U\implies M_1\abs{g(x)}\le\abs{f(x)}\le M_2\abs{g(x)}\); 
then set $M=M_2$, for $x\ge\max\{T,U\}$
\(M_2\abs{g(x)}\ge\abs{f(x)}>M_2\abs{g(x)},\); 
which can't happen. If $f\in o(g(x))\cap\omega(g(x))$,
\(\forall M>0\colon\exists T\in\bbR\colon x\ge T\implies\abs{f(x)}<M\abs{g(x)}\); 
\(\forall M>0\colon\exists U\in\bbR\colon x\ge U\implies\abs{f(x)}>M\abs{g(x)}\); 
then set any fixed $M$,
\(M\abs{g(x)}<\abs{f(x)}<M\abs{g(x)},\); 
which can't happen.

For any fixed $g(x)$,
\(o(g(x))\subseteq O(g(x))\); 

Proof:
Peak any $M>0$.

For any fixed $g(x)$,
\(\omega(g(x))\subseteq\Omega(g(x))\); 

Proof:
Peak any $M>0$.

For any fixed $g(x)$,
\(\Theta(g(x))=O(g(x))\cap\Omega(g(x))\); 

Proof:
Prove
\(\Theta(g(x))\subseteq O(g(x))\cap\Omega(g(x))\); 
\(\exists M>0,T\in\bbR\colon x\ge T\implies\abs{f(x)}\le M\abs{g(x)}\); 
\(\exists N>0,U\in\bbR\colon x\ge U\implies\abs{f(x)}\ge N\abs{g(x)}\); 
then
\(\exists M,N>0,\max\{T,U\}\in\bbR\colon x\ge \max\{T,U\}\implies N\abs{g(x)}\le\abs{f(x)}\le M\abs{g(x)}\); 

Prove
\(\Theta(g(x))\supseteq O(g(x))\cap\Omega(g(x))\); 
\(\exists M_1,M_2>0,T\in\bbR\colon x\ge T\implies M_1\abs{g(x)}\le\abs{f(x)}\le M_2\abs{g(x)}\); 
then
\(\exists M_2>0,T\in\bbR\colon x\ge T\implies\abs{f(x)}\le M_2\abs{g(x)}\); 
\(\exists M_1>0,T\in\bbR\colon x\ge T\implies\abs{f(x)}\ge M_1\abs{g(x)}\); 

For any fixed $g(x)$,
\(O(g(x))\setminus (o(g(x))\cup\Theta(g(x)))\ne\varnothing\); 

Proof:
Counterexample:
\(f(x)=\bcs 1,\quad & x\in\bbQ\\0,\quad & x\notin\bbQ\ecs\); 
\(g(x)=1\); 
$f\in O(g(x))$:
\(\exists 1>0,1\in\bbR\colon x\ge 1\implies\abs{f(x)}\le \abs{g(x)}\); 
$f\notin o(g(x))$:
\(\forall 1>0,T\in\bbR\colon\exists x\ge T,x\in\bbQ\colon\abs{f(x)}=\abs{g(x)}\); 
$f\notin\Theta(g(x))$:
\(\forall M_1>0,x\in\bbR,x\notin\bbQ\colon M_1\abs{g(x)}=M_1>\abs{f(x)}=0\); 

For any fixed $g(x)$,
\(\Omega(g(x))\setminus (\omega(g(x))\cup\Theta(g(x)))\ne\varnothing\); 

Proof:
Counterexample:
\(f(x)=1\); 
\(g(x)=\bcs 1,\quad & x\in\bbQ\\0,\quad & x\notin\bbQ\ecs\); 
$f\in \Omega(g(x))$:
\(\exists 1>0,1\in\bbR\colon x\ge 1\implies\abs{f(x)}\ge \abs{g(x)}\); 
$f\notin \omega(g(x))$:
\(\forall 1>0,T\in\bbR\colon\exists x\ge T,x\in\bbQ\colon\abs{f(x)}=\abs{g(x)}\); 
$f\notin\Theta(g(x))$:
\(\forall M_2>0,x\in\bbR,x\notin\bbQ\colon M_2\abs{g(x)}=0<\abs{f(x)}=1\); 

\(f(x)\in O(g(x))\iff g(x)\in\Omega(f(x))\); 

Proof:
$\implies$:
\(\exists M>0,T\in\bbR\colon x\ge T\implies\abs{f(x)}\le M\abs{g(x)}\); 
then
\(\exists \frac{1}{M}>0,T\in\bbR\colon x\ge T\implies\abs{g(x)}\ge \frac{1}{M}\abs{g(x)}\); 
$\Leftarrow$:
\(\exists M>0,T\in\bbR\colon x\ge T\implies\abs{f(x)}\ge M\abs{g(x)}\); 
then
\(\exists \frac{1}{M}>0,T\in\bbR\colon x\ge T\implies\abs{g(x)}\le \frac{1}{M}\abs{g(x)}\); 

\(f(x)\in o(g(x))\iff g(x)\in\omega(f(x))\); 

Proof:
$\implies$:
\(\forall M>0\colon\exists T\in\bbR\colon x\ge T\implies\abs{f(x)}<M\abs{g(x)}\); 
then
\(\forall \frac{1}{M}>0\colon\exists T\in\bbR\colon x\ge T\implies\abs{g(x)}>\frac{1}{M}\abs{f(x)}\); 
$\Leftarrow$:
\(\forall M>0\colon\exists T\in\bbR\colon x\ge T\implies\abs{f(x)}>M\abs{g(x)}\); 
then
\(\forall \frac{1}{M}>0\colon\exists T\in\bbR\colon x\ge T\implies\abs{g(x)}<\frac{1}{M}\abs{f(x)}\); 

For any fixed $g(x)$,
\(\exists f(x)\colon f(x)\notin O(g(x))\cup\Omega(g(x))\); 

Proof:
\(f(x)=\bcs 1,\quad & x\in\bbQ\\x,\quad & x\notin\bbQ\ecs\); 
\(g(x)=\bcs x,\quad & x\in\bbQ\\1,\quad & x\notin\bbQ\ecs\); 
For $x\in \bbR\setminus\bbQ$,
\(\nexists M>0,T\in\bbR\colon x\ge T\implies\abs{f(x)}\le M\abs{g(x)}\); 
For $x\in \bbQ$,
\(\nexists M>0,T\in\bbR\colon x\ge T\implies\abs{f(x)}\ge M\abs{g(x)}\); 


\tb{Binary relation}

$\Theta(g(x))$ is an equivalence class.

Proof:
Transitive:
\(\exists M_1,M_2>0,T\in\bbR\colon x\ge T\implies M_1\abs{g(x)}\le\abs{f(x)}\le M_2\abs{g(x)}\); 
\(\exists N_1,N_2>0,U\in\bbR\colon x\ge U\implies N_1\abs{h(x)}\le\abs{g(x)}\le N_2\abs{h(x)}\); 
then
\(\exists M_1N_1,M_2N_2>0,\max\{T,U\}\in\bbR\colon x\ge \max\{T,U\}\implies M_1N_1\abs{h(x)}\le\abs{f(x)}\le M_2N_2\abs{h(x)}\); 
Reflexive:
\(f(x)\in\Theta(f(x))\); 
\(\exists 1,1>0,a\in\bbR\colon x\in\bbR\implies 1\abs{f(x)}\le\abs{f(x)}\le 1\abs{f(x)}\); 
Symmetric:
\(\exists M_1,M_2>0,T\in\bbR\colon x\ge T\implies M_1\abs{g(x)}\le\abs{f(x)}\le M_2\abs{g(x)}\); 
then
\(\exists \frac{1}{M_2},\frac{1}{M_1}>0,T\in\bbR\colon x\ge T\implies \frac{1}{M_2}\abs{f(x)}\le\abs{f(x)}\le\frac{1}{M_1}\abs{f(x)}\); 


$O(g(x))$ is partial-ordered set with $\Theta$ as $=$.

Proof:
Transitive:
\(\exists M>0,T\in\bbR\colon x\ge T\implies\abs{f(x)}\le M\abs{g(x)}\); 
\(\exists N>0,U\in\bbR\colon x\ge U\implies\abs{g(x)}\le N\abs{h(x)}\); 
then
\(\exists MN>0,\max\{T,U\}\in\bbR\colon x\ge\max\{T,U\}\implies\abs{f(x)}\le MN\abs{h(x)}\); 
Reflexive:
\(\exists M>0,T\in\bbR\colon x\ge T\implies\abs{f(x)}\le M\abs{g(x)}\); 
then
\(\exists \frac{1}{M}>0,T\in\bbR\colon x\ge T\implies\abs{g(x)}\le \frac{1}{M}\abs{f(x)}\); 
Antisymmetric:
\(\exists M>0,T\in\bbR\colon x\ge T\implies\abs{f(x)}\le M\abs{g(x)}\); 
\(\exists N>0,U\in\bbR\colon x\ge U\implies\abs{g(x)}\le N\abs{f(x)}\); 
then
\(\exists M,\frac{1}{N}>0,\max\{T,U\}\in\bbR\colon x\ge \max\{T,U\}\implies\frac{1}{N}\abs{g(x)}\le\abs{f(x)}\le M\abs{g(x)}\); 


$\Omega(g(x))$ is partial-ordered set with $\Theta$ as $=$.

Proof:
Similar to that of $O$.


$o(g(x))$ is a strictly partially-ordered set.

Proof:
Asymmetry:
\(f(x)\in o(g(x))\iff g(x)\in\omega(f(x))\implies g\notin o(f(x))\); 
Transitivity:
\(\forall M>0\colon\exists T\in\bbR\colon x\ge T\implies\abs{f(x)}<M\abs{g(x)}\); 
\(\exists U\in\bbR\colon x\ge U\implies\abs{g(x)}<\abs{h(x)}\); 
then
\(\forall M>0\colon\exists \max\{T,U\}\in\bbR\colon x\ge \max\{T,U\}\implies\abs{f(x)}<M\abs{h(x)}\); 


$\omega(g(x))$ is a strictly partially-ordered set.

Proof:
Similar to that of $o$.


\tb{Integral test (for convergence) or Maclaurin–Cauchy test}

Suppose $f(x)\colon[a,\infty)\to\bbR$ for some $a\in\bbN$ is monotone and $\lim_{x\to\infty}f(x)=0$. Then $\int_a^{\infty}f(x)\dd{x}$ is convergent iff $\sum_{k=a}^{\infty}f(k)$ is convergent. And,

If $\int_a^{\infty}f(x)\dd{x}$ is convergent and $f(x)$ is monotone decreasing, for any $n\in\bbN_{\ge a}$,
\(\int_n^{\infty}f(x)\dd{x}\le\sum_{k=n}^{\infty}f(k)\le f(n)+\int_n^{\infty}f(x)\dd{x},\); 
called remainder estimate for integral test. Note that for $n\in\bbN_{>a}$, $f(n)+\int_n^{\infty}f(x)\dd{x}\le\int_{n-1}^{\infty}f(x)\dd{x}$.

If $\int_a^{\infty}f(x)\dd{x}$ is convergent and $f(x)$ is monotone increasing, for any $n\in\bbN_{\ge a}$,
\(\int_n^{\infty}f(x)\dd{x}\ge\sum_{k=n}^{\infty}f(k)\ge f(n)+\int_n^{\infty}f(x)\dd{x},\); 
called remainder estimate for integral test. Note that for $n\in\bbN_{>a}$, $f(n)+\int_n^{\infty}f(x)\dd{x}\ge\int_{n-1}^{\infty}f(x)\dd{x}$.

Proof:
Consider the case where $f(x)$ is monotone increasing. The case where $f(x)$ is monotone decreasing can be proven similarly. First consider the case where $\int_a^{\infty}f(x)\dd{x}$ is convergent. This implies that $f(x)$ is non-positive. Thus, for any fixed $n\in\bbN_{\ge a}$,
\(\int_n^{n+1}f(x)\dd{x}\ge f(n),\); 
and for any $k\geq n+1$,
\(\int_k^{k+1}f(x)\dd{x}\ge f(k)\ge\int_{k-1}^kf(x)\dd{x}.\); 
For any $N\in\bbN_{\ge n}$, summing over all $k\in\bbN_{\ge n}$, we obtain, 
\(\int_n^{N+1}f(x)\dd{x}\ge\sum_{k=n}^Nf(k)\ge f(n)+\int_n^Nf(x)\dd{x}.\); 
Take limit $\lim_{N\to\infty}$, we obtain
\(\int_n^{\infty}f(x)\dd{x}\ge\sum_{k=n}^{\infty}f(n)\ge f(n)+\int_n^{\infty}f(x)\dd{x}.\); 
Second, consider the case where $\int_a^{\infty}f(x)\dd{x}$ is divergent. Since $f(x)$ is monotone. Either $\int_a^{\infty}f(x)\dd{x}=-\infty$ or $\int_a^{\infty}f(x)\dd{x}=\infty$. For the former, the supremum of $f(x)$ is non-positive, and $\sum_{k=a}^{\infty}f(k)=-\infty$. For the latter, there exists $x_0>a$ such that $f(x)>0$ for all $x\ge x_0$, and thus $\sum_{k=a}^{\infty}f(k)=\infty$.


\tb{p-test (for improper integral)}

For any fixed $a>0$, $\int_a^{\infty}\frac{1}{x^p}\dd{x}$ converges iff $p>1$, and if so,
\(\int_a^{\infty}\frac{1}{x^p}\dd{x}=\frac{1}{\qty(p-1)a^{p-1}};\); 
$\int_0^a\frac{1}{x^p}\dd{x}$ converges iff $p<1$, and if so,
\(\int_0^a\frac{1}{x^p}\dd{x}=-\frac{1}{\qty(p-1)a^{p-1}}.\); 
Proof:
\(\int\frac{1}{x^p}\dd{x}=\begin{cases}
-\frac{1}{\qty(p-1)x^{p-1}},\quad&p\neq 1\\
\ln\qty(x),\quad&p=1
\end{cases}\); 
\(\forall p<1\colon\nexists\lim_{x\to\infty}-\frac{1}{\qty(p-1)x^{p-1}}\); 
\(\nexists\lim_{x\to\infty}\ln\qty(x)\); 
\(\forall p>1\colon\lim_{x\to\infty}-\frac{1}{\qty(p-1)x^{p-1}}=0\quad\int_a^{\infty}\frac{1}{x^p}\dd{x}=\frac{1}{\qty(p-1)a^{p-1}}\); 
\(\forall p>1\colon\nexists\lim_{x\to 0^+}-\frac{1}{\qty(p-1)x^{p-1}}\); 
\(\nexists\lim_{x\to 0^+}\ln\qty(x)\); 
\(\forall p<1\colon\lim_{x\to 0^+}-\frac{1}{\qty(p-1)x^{p-1}}=0\quad\int_0^a\frac{1}{x^p}\dd{x}=-\frac{1}{\qty(p-1)a^{p-1}}\); 

\tb{Generalized p-test (for improper integral)}

Define iterated logarithmic function $L_n\qty(x)$ as
\(L_{0}\qty(x)=x,\quad x\in\bbR_{>0}\); 
\(L_{n}\qty(x)=\ln\qty(L_{n-1}\qty(x)),\quad n\in\bbN,x\in\bbR_{>0}.\); 
For any fixed $a>1,k\in\bbN$,
\(\int_a^{\infty}\frac{1}{\qty(\prod_{i=0}^{k-1}L_{i}\qty(x))\qty(L_{k}\qty(x))^m}\dd{x}\); 
converges iff $m>1$.

Proof:
Let
\(u=L_k\qty(x)\); 
\(\dd{u}=\frac{\dd{x}}{\prod_{i=0}^{k-1}L_{i}\qty(x)}\); 
\(\int\frac{1}{\qty(\prod_{i=0}^{k-1}L_{i}\qty(x))\qty(L_{k}\qty(x))^m}\dd{x}=\int\frac{1}{u^m}\dd{u}\); 
By p-test, the statement holds.

For any fixed $0<a<1$,
\(\int_0^a\frac{1}{x(\ln x)^m}\dd{x}\); 
converges iff $m<1$.

Proof:
By substitution of $u=\ln x$.


\tb{Generalized p-test for series or generalized p-series test}

Define iterated logarithmic function $L_n\qty(x)$ as
\(L_{0}\qty(x)=x,\quad x\in\bbR_{>0}\); 
\(L_{n}\qty(x)=\ln\qty(L_{n-1}\qty(x)),\quad n\in\bbN,x\in\bbR_{>0}.\); 
For any fixed $a\in\bbN$,
\(\sum_{n=a}^{\infty}\frac{1}{\qty(\prod_{i=0}^{k-1}L_{i}\qty(n))\qty(L_{k}\qty(n))^m}\); 
converges iff $m>1$.

Proof:
By integral test and generalized p-test for improper integral, the statement holds.


\tb{Generalized Master Theorem}

For a positive valued net $T\qty(n)$ defined on $n=b^h$ for any $h\in\bbN_0$ as
\(T\qty(n)=aT\qty(\frac{n}{b})+f\qty(n)\); 
for some constants $a>0$ and $b>1$, and non-negative-valued function $f\qty(n)$,

If for some $c<\log_ba$,
\(f\qty(n)=O\qty(n^c),\); 
then
\(T\qty(n)=\Theta\qty(n^{\log_ba}).\); 

Define iterated logarithmic function $L_n\qty(x)$ as
\(L_{0}\qty(x)=x,\quad x\in\bbR_{>0}\); 
\(L_{n}\qty(x)=\ln\qty(L_{n-1}\qty(x)),\quad n\in\bbN,x\in\bbR_{>0}.\); 

If for some positive integers $k$ and positive number $m$, 
\(f\qty(n)=\Theta\qty(n^{\log_ba}\qty(\prod_{i=1}^{k-1}\qty(L_{i}\qty(n))^{-1})\qty(L_{k}\qty(n))^{m-1}),\); 
then
\(T\qty(n)=\Theta\qty(n^{\log_ba}\qty(L_{k}\qty(n))^m).\); 

If for some positive integer $k$, 
\(f\qty(n)=\Theta\qty(n^{\log_ba}\prod_{i=1}^k\qty(L_{i}\qty(n))^{-1}),\); 
then
\(T\qty(n)=\Theta\qty(n^{\log_ba}L_{k-1}\qty(n)).\); 

If for some positive integers $k$ and some positive number $m$,
\(f\qty(n)=\Theta\qty(n^{\log_ba}\qty(\prod_{i=1}^{k-1}\qty(L_{i}\qty(n))^{-1})\qty(L_{k}\qty(n))^{-m-1}),\); 
then
\(T\qty(n)=\Theta\qty(n^{\log_ba}).\); 

If $f\qty(n)=\Omega\qty(n^c)$ for some $c>\log_ba$, and the regularity condition that $\exists m\in\bbR,k<1$ such that $n\ge m$ implies
\(af\qty(\frac{n}{b})\le kf\qty(n)\); 
is satisfied, then
\(T\qty(n)=\Theta\qty(f\qty(n)).\); 

Proof:
Let $s=\log_ba$, $n=b^h$, and $f\qty(x)=x^sg\qty(x)$. At level $j$, there are $a^j$ nodes, each of size $\frac{n}{b^j}$. Thus, the total non-recursive work at level $j$ is
\(a^jf\qty(\frac{n}{b^j}).\); 
Hence
\(
T\qty(n)=\sum_{j=0}^{h-1}a^jf\qty(\frac{n}{b^j})+n^sT\qty(1)
=\sum_{j=0}^{h-1}b^{sj}f\qty(\frac{b^h}{b^j})+\Theta\qty(n^s)
=n^s\sum_{j=0}^{h-1}g\qty(b^{h-j})+\Theta\qty(n^s)
=n^s\sum_{r=1}^hg\qty(b^r)+\Theta\qty(n^s)
\); 
Case 1: If for some $c<\log_ba$,
\(f\qty(n)=O\qty(n^c),\); 
then
\(T\qty(n)=\Theta\qty(n^{\log_ba}).\); 
Case 2:
\(L_{k}\qty(b^x)=\Theta\qty(L_{k-1}\qty(x))\); 
\(f\qty(n)=\Theta\qty(n^{\log_ba}\qty(\prod_{i=1}^{k-1}\qty(L_{i}\qty(n))^{-1})\qty(L_{k}\qty(n))^{m-1})\); 
\(g\qty(b^r)=\Theta\qty(\qty(\prod_{i=1}^{k-2}\qty(L_{i}\qty(r))^{-1})\qty(L_{k-1}\qty(r))^{m-1})\); 
By generalized p-test for series: Sub-case 1 $m>0$:
\(T\qty(n)=\Theta\qty(n^{\log_ba}\qty(L_{k}\qty(n))^m).\); 
Sub-case 2 $m=0$:
\(T\qty(n)=\Theta\qty(n^{\log_ba}L_{k-1}\qty(n)).\); 
Sub-case 3 $m=0$:
\(T\qty(n)=\Theta\qty(n^{\log_ba}).\); 
Case 3: If $f\qty(n)=\Omega\qty(n^c)$ for some $c>\log_ba$, and the regularity condition that $\exists m\in\bbR,0<k<1$ such that $n\ge m$ implies
\(af\qty(\frac{n}{b})\le kf\qty(n)\); 
is satisfied, then
\(af\qty(\frac{n}{b})\le kf\qty(n)\); 
\(a^jf\qty(\frac{n}{b^j})\le a^{j-1}kf\qty(\frac{n}{b^{j-1}})\le k^jf\qty(n)\); 
\(
T\qty(n)=\sum_{j=0}^{h-1}a^jf\qty(\frac{n}{b^j})+n^sT\qty(1)
=f\qty(n)+\sum_{j=1}^{h-1}O\qty(k^jf\qty(n))+\Theta\qty(n^s)
\); 
\(\sum_{j=1}^{\infty}k^jf\qty(n)=\frac{f\qty(n)}{1-k}\); 
\(T\qty(n)=f\qty(n)+O\qty(f\qty(n))+\Theta\qty(n^s)=\Theta\qty(f\qty(n))\); 

Terms:

$h$ is called recursion level.

$\log_ba$ is called critical exponent $c_{crit}$.

$T\qty(1)$ is called leaf cost.

At each node, the first term is called recursive work, and the second term is called non-recursive work.

For original master theorem, the second case is

If for some positive integers $m$, 
\(f\qty(n)=\Theta\qty(n^{\log_ba}\qty(\ln n)^{m-1}),\); 
then
\(T\qty(n)=\Theta\qty(n^{\log_ba}\qty(\ln n)^m).\); 

If
\(f\qty(n)=o\qty(n^{\log_ba}\qty(\ln n)^{-1}),\); 
then
\(T\qty(n)=\Theta\qty(n^{\log_ba}\ln\qty(\ln n)).\); 

If for some positive integers $m$,
\(f\qty(n)=\Theta\qty(n^{\log_ba}\qty(\ln n)^{-m-1}),\); 
then
\(T\qty(n)=\Theta\qty(n^{\log_ba}).\); 

\tb{Stirling's approximation}

\(n!=\sqrt{2\pi n}\qty(\frac{n}{e})^n\qty(1+O\qty(\frac{1}{n}))\);

\tb{特徵方程式}

定義：對一個給定$k$階常係數線性遞迴數列：
\(<a_n>: a_n=\sum_{i=1}^k c_i \cdot a_{n-i} + f(n)\); 
，其特徵方程式為$x^k-\sum_{i=1}^k c_i \cdot x^{k-i}=0$，特證方程式的根稱其特徵根。

若一個給定$k$階常係數線性齊次遞迴數列$<a_n>$有$a$相異特徵根$q_1, q_2, q_3, \ldots, q_a$，其中第$i$個根（$1\leq i\leq a$）為$m_i$重根，則$<a_n>$之一般式為：
\(a_n=\sum_{i=1}^{a} \left( \sum_{j=1}^{m_i} \left(A_{i,j}\cdot n^{j-1}\right) \cdot q_i^{\phantom{i}n} \right)\); 

若給定一個$k$階常係數線性非齊次遞迴數列：
\(<a_n>: a_n=\sum_{i=1}^k c_i \cdot a_{n-i} + f(n)\); 
則令：
\(a_n = a'_n+a''_n \); 
使得$<a'_n>$為$k$階常係數線性齊次遞迴數列，$<a''_n>$為$k$階常係數線性非齊次遞迴數列且一次項同$<a'_n>$且$f(n)$項同$<a_n>$。
\(a''_n = g(n),\quad\deg(g(n))\leq \deg(f(n))\); 
解出$k$階常係數線性非齊次遞迴數列$<a'_n>$與$g(n)$ ，即可解出$<a_n>$。

\tb{Linear Data Structures}

An unordered list is a list where each element holds a relative position with respect to the others. Each element is called node. If applicable, the first node is called head and the last node is called tail.
An ordered list is a list which are ordered based upon some underlying characteristic of the elements.
A stack is a last-in-first-out (LIFO) list. Adding an element is called push and is O(1), the last pushed element is called top, and the first pushed element is called bottom. Removing and returning top is called pop and is O(1). Peeking top is called top and is O(1).
A queue is first-in-first-out (FIFO) list. Adding an element is called enqueue and is O(1), the first enqueued element is called front, and the last enqueued element is called back. Removing and returning front is called dequeue and is O(1). Peeking front is called peek and is O(1).
A doubly-ended queue (deque) is a list that is simultaneously a stack and a queue, that is, has O(1) peeking, pushing, and popping from both ends.
A dynamic array is a contiguous array that reserves a new capacity elsewhere and moves to it whenever its size reaches its capacity. It has O(1) random access and popping back, amortized O(1) pushing back, but O(n) removal and insertion. It is a stack with front as bottom and back as top.
A singly linked list is an unordered list where each node stores only the next node. It has O(1) peeking, pushing, and popping front and insertion after node, but O(n) random access, removal of node, and insertion before node. If the tail is stored, it has O(1) peeking and pushing back. It is a stack with head as top and tail as bottom. If the tail is stored, it is a queue with the tail as front and head as back.
A doubly linked list is an unordered list where each node stores the previous node and the next node. It has O(1) peeking, pushing, and popping front and back, insertion after and before node, and removal of node, but O(n) random access. It is a deque.
An implementation of deque is a dynamic array of fixed-size blocks, which has O(1) random access and popping front and back, amortized O(1) pushing front and back, but O(n) removal and insertion.
To keep track of the minimum element of a stack in O(1) time with O(n) extra space. One can maintain the original stack of elements and an extra stack of extrema. When an element is pushed to the original stack, it is compared with the top of the extrema stack and the minimum one is pushed to the extrema stack (for the first element, it is pushed to the extrema stack). When an element is popped from the original stack, an element is popped from the extrema stack. This way, the top of the extrema stack is always the minimum of the elements in the original stack.

\lstinputlisting[language=Python]{basic.py}
\lstinputlisting[language=Python]{str.py}
\lstinputlisting[language=Python]{list.py}
\lstinputlisting[language=Python]{tuple.py}
\lstinputlisting[language=Python]{set.py}
\lstinputlisting[language=Python]{dict.py}
\lstinputlisting[language=Python]{SinglyLinkedList.py}
\lstinputlisting[language=Python]{DoublyLinkedList.py}
\lstinputlisting[language=Python]{BlockListDeque.py}
\end{multicols}
\end{document}
